% Subsection 2.2.2: LoRA 原理与优势

LoRA（Low-Rank Adaptation）是目前应用最广泛、稳定性最强的参数高效微调方法之一\cite{hu2021lora}，其核心思想是在不改变基座模型权重的前提下，在 Transformer 的注意力模块中引入低秩分解矩阵，仅通过训练这些少量参数实现任务适配。

在结构上，LoRA 将下游任务带来的参数更新量 $\Delta W$ 约束在低维空间中：对于原始权重矩阵 $W \in \mathbb{R}^{d \times k}$，更新量被显式分解为两个低秩矩阵的乘积
\begin{equation}
    W' = W + \Delta W = W + BA,\quad A \in \mathbb{R}^{r \times k},\ B \in \mathbb{R}^{d \times r},\ r \ll \min(d, k),
\end{equation}
其中秩 $r$ 通常取 $8$ 至 $64$。训练时冻结 $W$，仅优化 $A$ 与 $B$；前向传播时，原始注意力输出与低秩分支输出叠加得到最终表示，既保留了模型原本的代码理解与生成能力，又能针对性地学习算法题的解题模式与逻辑规律。低秩分支与冻结主干的并联关系如图~\ref{fig:ch2-lora-structure} 所示：输入 $x$ 同时进入冻结主干 $W$ 与可训练旁路 $A \to B$，二者输出相加得到最终隐藏态 $h = Wx + BAx$，旁路的低秩约束 $r \ll \min(d, k)$ 决定了可训练参数量仅为全量微调的极小比例。

\begin{figure}[!htb]
  \centering
  \begin{tikzpicture}[
    node distance=5mm and 6mm,
    inout/.style={draw, circle, minimum size=7mm, font=\footnotesize, inner sep=0pt},
    frozen/.style={draw, fill=black!10, rounded corners=2pt, align=center, minimum height=10mm, minimum width=18mm, font=\footnotesize},
    train/.style={draw, fill=blue!10, rounded corners=2pt, align=center, minimum height=8mm, minimum width=14mm, font=\footnotesize},
    sum/.style={draw, circle, minimum size=6mm, font=\footnotesize, inner sep=0pt},
    arr/.style={-{Stealth[length=2mm]}, thick},
    lbl/.style={font=\scriptsize, align=center},
  ]
    \node[inout] (x) {$x$};
    \node[frozen, right=10mm of x] (W) {$W$\\ \scriptsize 冻结};
    \node[train, above=8mm of W, xshift=14mm] (A) {$A \in \mathbb{R}^{r\times k}$};
    \node[train, right=of A] (B) {$B \in \mathbb{R}^{d\times r}$};
    \node[sum, right=20mm of W] (s) {$+$};
    \node[inout, right=of s] (h) {$h$};

    \draw[arr] (x) -- (W);
    \draw[arr] (x.east) to [bend left=35] (A.west);
    \draw[arr] (A) -- node[midway, above, lbl]{$r$ 维} (B);
    \draw[arr] (B.east) to [bend left=20] (s.north);
    \draw[arr] (W) -- (s);
    \draw[arr] (s) -- (h);

    \node[lbl, below=1mm of W] {主干分支：$Wx$};
    \node[lbl, above=1mm of A.north -| s] {可训练低秩旁路：$BAx$};
    \node[lbl, below=2mm of h, xshift=-6mm] {$h = Wx + BAx$};
  \end{tikzpicture}
  \caption{LoRA 在单层线性变换上的结构：冻结主干 $W$ 与可训练低秩旁路 $A\!\to\!B$ 并联，输出叠加}
  \label{fig:ch2-lora-structure}
\end{figure}

相比于全量微调和其他轻量适配方法，LoRA 具有多重优势：
\begin{itemize}
    \item \textbf{参数量极小}：可训练参数通常不到总参数的 $1\%$，显著降低显存占用与训练成本；
    \item \textbf{零推理开销}：推理阶段可将低秩矩阵与原权重合并，不引入额外计算开销，部署友好；
    \item \textbf{训练稳定}：在小规模数据上不易过拟合，更适合算法代码这类对逻辑正确性要求严格的场景；
    \item \textbf{适配器可插拔}：支持多适配器独立存储与快速切换，为开展对照实验提供了极大便利。
\end{itemize}
凭借这些特点，LoRA 成为本文实现模型高效适配的核心技术支撑。